\setlength{\footskip}{8mm}

\chapter{CONCLUSION}

This thesis set out from a simple observation about human communication: in therapy, what a client says and how they say it do not always agree. Text-only AI counselors are blind to this gap. They read the words ``I am fine'' and have no way to hear the tremor underneath them. We proposed that emotional dissonance, computed as the Valence-Arousal mismatch between a client's text and their voice, can serve as an explicit signal that guides a CBT-oriented LLM therapist toward deeper, more empathetic, and more clinically appropriate responses.

We implemented this idea as a complete multimodal pipeline. Client utterances are analyzed in two parallel channels: vad-bert extracts text-based Valence-Arousal values, and WavLM extracts speech-based Valence-Arousal values from the vocal delivery, with librosa providing natural-language vocal descriptors. The per-dimension delta between the two channels, flagged when either exceeds a fixed threshold of 0.5, is injected directly into the therapist's prompt as a dissonance signal, instructing the model to gently explore what might lie beneath the surface when words and voice diverge.

Our experiments answer both research questions affirmatively. For RQ1, across 100 controlled CBT dialogues on GPT-4o mini, the Dissonance-Aware condition outperformed both text-only baselines (Baseline and Emotion) on all seven therapeutic quality metrics judged by a blind GPT-4o evaluator, with an average uplift of +0.64 points and the largest gains on the CTRS dimensions that measure collaboration (+0.80), guided discovery (+0.68), strategy (+0.62), and focus (+0.59); against the Multimodal (Vocal-Aware) condition it remained near-equivalent, leading most clearly on CTRS Focus. For RQ2, the advantage replicated across all four tested LLM backbones (GPT-4o mini, Claude Sonnet 4.6, Qwen-2.5-72B-Instruct, and DeepSeek V3.2), spanning dense and Mixture-of-Experts architectures and both closed and open-weight models. The magnitude of the gain scaled inversely with baseline capability: the weakest backbone gained the most, and GPT-4o mini with the dissonance signal moved substantially closer to Claude Sonnet 4.6's unaided baseline, at roughly a twentieth of the inference cost.

The contributions of this work are fourfold. First, we introduced a multimodal framework that aligns independently extracted text and speech emotional representations for therapist response generation. Second, we operationalized emotional dissonance as a prompt-level therapeutic signal, reframing what prior work treated as multimodal failure into a clinically useful cue of hidden affect. Third, we established a controlled four-condition experimental design, paired on shared problem seeds, that isolates the contribution of each emotional metadata type. Fourth, we delivered a reproducible blind evaluation protocol and demonstrated the framework's robustness across four architecturally diverse LLM backbones.

The broader implication is encouraging for accessible mental health technology: emotionally perceptive therapeutic AI does not require frontier-scale models or expensive training. A modest model, told explicitly when a client's voice disagrees with their words, can approach the unaided performance of a model many times its price. For real deployments where cost and privacy constrain model choice, this is a practical path toward therapists that hear not only what is said, but how it is said.

We close with candor about what remains. Our evidence comes from synthetic, English-language dialogues judged by an LLM, and our TTS client voices are produced by an alpha-stage engine. The immediate next steps are clear: cross-turn dissonance tracking with an adaptive threshold, and validation with real human voices and clinical expert raters. We regard the present results as a foundation: emotional dissonance is measurable, injectable, and demonstrably useful, and the door is now open to richer, more human-like multimodal therapeutic reasoning.
